% Einheit 2 von 2 – Vorgegebene Regeln anwenden (bank/integrationsregeln/e2.jsonl, zusammenbau v0.1)
%% TODO Zeitmarke Einheit 2: Marken-Zeile nicht lesbar
%% TODO Zweigzeile Teil 1: was hier gelernt wird (Satz in Schülersprache aus dem Einheitstitel) – Einheit 2
\einheitenkopf[e2]{Einheit 2 von 2 · Vorgegebene Regeln anwenden}
\zweigzeile{Abitur LK}
\setcounter{aufgabe}{6}

%% TODO Titel: Ich-kann-Satz fehlt in der Bank (Platzhalter: Kettenname) – Nr. 7
%% TODO Anweisung: ein Satz über der ersten Teilaufgabe fehlt in der Bank – Nr. 7
\begin{aufgabe}{Vorgegebene Regeln anwenden}
%% TODO \rechenplatz{n}: Schrittzahl des Grundfalls fehlt in der Bank (nur bei fester Schreibform, 2.3 b)
\begin{teile}
\teil $2x \cdot e^{x^2 + 1}$: Welcher Teil ist g?\\ \kreuz{der Exponent}\\ \kreuz{der Vorfaktor}\\ \kreuz{die ganze e-Potenz}
\teil $f(x) = 2x \cdot e^{x^2 + 3}$. g und g'? $g(x) = $ \leerfeld, $g'(x) = $ \leerfeld
\teil Schreibe $f(x) = 2x \cdot e^{1 - x^2}$ mit $g'(x) \cdot e^{g(x)}$. \leerfeld
\teil Gegeben ist die Regel: Das Integral von u bis v über $g'(x) \cdot e^{g(x)}$ ist $\left[e^{g(x)}\right]_u^v$. Berechne das Integral von 0 bis 1 über $2x \cdot e^{x^2}$. \leerfeld
\teil Gegeben ist die Regel: Das Integral von u bis v über $g'(x) \cdot e^{g(x)}$ ist $\left[e^{g(x)}\right]_u^v$. Berechne das Integral von 0 bis 1 über $2x \cdot e^{2 - x^2}$ mithilfe dieser Regel. \hfill (Abitur 2022 LK) \\ \leerfeld
\end{teile}
\end{aufgabe}

%% TODO Titel: Ich-kann-Satz fehlt in der Bank (Platzhalter: Kettenname) – Nr. 8
%% TODO Anweisung: ein Satz über der ersten Teilaufgabe fehlt in der Bank – Nr. 8
\begin{aufgabe}{Vorgegebene Regeln anwenden – Fehler finden, Begründen}
\begin{teile}
\teil Gegeben ist die Regel: Das Integral von u bis v über $g'(x) \cdot e^{g(x)}$ ist $\left[e^{g(x)}\right]_u^v$. Für das Integral von 0 bis 1 über $2x \cdot e^{-x^2}$ rechnet Mia $\left[e^{-x^2}\right]_0^1$. Finde den Fehler.
\teil Begründe: Die Regel „das Integral über $g'(x) \cdot e^{g(x)}$ ist $\left[e^{g(x)}\right]$“ gilt nur für Integranden der Form $g' \cdot e^g$.
\end{teile}
\end{aufgabe}
